\documentclass[11pt]{article}

\usepackage[T1]{fontenc}
\usepackage{lmodern}
\usepackage{amsmath,amssymb,amsthm}
\usepackage[margin=1in]{geometry}
\usepackage{enumitem}
\usepackage[colorlinks=true,linkcolor=blue,citecolor=blue,urlcolor=blue]{hyperref}


% ---------- theorem-like environments ----------
\newtheorem{theorem}{Theorem}
\newtheorem{proposition}[theorem]{Proposition}
\newtheorem{lemma}[theorem]{Lemma}
\newtheorem{corollary}[theorem]{Corollary}
\theoremstyle{definition}
\newtheorem{definition}[theorem]{Definition}
\newtheorem{example}[theorem]{Example}
\theoremstyle{remark}
\newtheorem{remark}[theorem]{Remark}

% ---------- macros ----------
\newcommand{\R}{\mathbb{R}}
\newcommand{\C}{\mathbb{C}}
\newcommand{\one}{\mathbf{1}}
\newcommand{\T}{\top}
\renewcommand{\Re}{\operatorname{Re}}
\newcommand{\diag}{\operatorname{diag}}
\newcommand{\cbar}{\bar{c}}

\setlist[enumerate,1]{label=\textup{(\alph*)}, leftmargin=2.2em}

\title{Convergence of the normalized state vector \\ in the positive linear system $\dot{x} = Mx$}
\author{}
\date{}

\begin{document}
\maketitle

\noindent
Let $M$ be a Metzler matrix and let $x(t)$ solve $\dot{x} = Mx$ from a nonnegative, nonzero
initial state. These notes prove that the normalized solution $x(t)/\sum_i x_i(t)$ converges to a
nonnegative eigenvector of $M$, with no irreducibility assumption
(Theorem~\ref{thm:main}).

% =====================================================================
\section{Setting and main result}\label{sec:setting}
% =====================================================================

\paragraph{Notation.}
For a real vector or matrix $A$ we write $A \ge 0$ if every entry of $A$ is nonnegative, and
$A \gg 0$ if every entry is positive. $I$ is the identity matrix, $\one$ is the column vector of
ones (of the appropriate dimension), and $e_j$ is the $j$th standard basis vector. For an index set $R \subseteq \{1,\dots,n\}$,
$x_R = (x_i)_{i \in R}$ is the corresponding subvector of $x$, and $M_{RR} = (M_{ij})_{i,j \in R}$
is the corresponding principal submatrix of $M$; in both, the indices of $R$ are kept in ascending
order. Eigenvalues may be complex, and $\Re\lambda$ is
the real part of $\lambda$.

\paragraph{Standing hypotheses.}
$M$ is a real $n \times n$ matrix with $M_{ij} \ge 0$ whenever $i \ne j$ (a \emph{Metzler} matrix).
The initial state $x(0) \in \R^n$ satisfies $x(0) \ge 0$ and $x(0) \ne 0$, and $x(t)$, $t \ge 0$,
is the solution of $\dot{x} = Mx$. We write
\[
  W(t) = \sum_{i=1}^n x_i(t) = \one^\T x(t),
  \qquad
  p(t) = \frac{x(t)}{W(t)} .
\]
If $x_i$ is the abundance of type $i$, then $W$ is the total population and $p$ is the vector of
type frequencies.

\begin{definition}[Reachability]\label{def:reach}
There is an \emph{arc} from $j$ to $i$ if $i \ne j$ and $M_{ij} > 0$. Index $i$ is
\emph{reachable from} $j$ if $i = j$, or if there are indices $j = i_0, i_1, \dots, i_r = i$ with
an arc from $i_{s-1}$ to $i_s$ for each $s = 1,\dots,r$. We set
\begin{align*}
  S &= \{\, j : x_j(0) > 0 \,\} && \text{(the support of the initial state),}\\
  R &= \{\, i : i \text{ is reachable from some } j \in S \,\} && \text{(the reachable set; } S \subseteq R\text{),}\\
  B &= M_{RR} && \text{(the matrix of the reduced system).}
\end{align*}
\end{definition}

An arc from $j$ to $i$ means that type $j$ contributes directly to the growth of type $i$.

\begin{theorem}\label{thm:main}
Under the standing hypotheses:
\begin{enumerate}
  \item\label{thm:main:support} For every $t > 0$, $x_i(t) > 0$ if $i \in R$; and $x_i(t) = 0$ for
        all $t \ge 0$ if $i \notin R$. In particular $W(t) > 0$ for all $t \ge 0$, so $p(t)$ is
        well defined.
  \item\label{thm:main:mu} The number
        $\mu = \max\{\, \Re\lambda : \lambda \text{ is an eigenvalue of } B \,\}$
        is itself an eigenvalue of $B$; in particular it is real.
  \item\label{thm:main:limit} The limit $p^* = \lim_{t \to \infty} p(t)$ exists, and
        \[
          p^* \ge 0, \qquad \one^\T p^* = 1, \qquad p^*_i = 0 \ \text{ for } i \notin R,
          \qquad M p^* = \mu\, p^* .
        \]
        In particular $\mu$ is an eigenvalue of $M$, with the nonnegative eigenvector $p^*$.
\end{enumerate}
\end{theorem}

The eigenvalue $\mu$ belongs to the reduced matrix $B$. It need not be the eigenvalue of $M$ with
the largest real part (Example~\ref{ex:reduction}).

% =====================================================================
\section{Background}\label{sec:background}
% =====================================================================

The proof uses three standard facts.

\begin{lemma}[The matrix exponential]\label{lem:exp}
For a real or complex square matrix $A$ and real $t$, the series
$e^{At} = \sum_{k=0}^{\infty} \frac{t^k}{k!} A^k$ converges, and:
\begin{enumerate}
  \item\label{lem:exp:ode} $\frac{d}{dt} e^{At} = A e^{At}$ and $e^{A \cdot 0} = I$. Hence the
        solution of $\dot{x} = Ax$ with initial value $x(0)$ is $x(t) = e^{At} x(0)$.
  \item\label{lem:exp:semigroup} $e^{A(s+t)} = e^{As} e^{At}$ for all real $s$, $t$.
  \item\label{lem:exp:shift} $e^{(A + cI)t} = e^{ct} e^{At}$ for every scalar $c \in \C$.
  \item\label{lem:exp:eig} If $Au = \lambda u$ with $\lambda \in \C$, then $e^{At} u = e^{\lambda t} u$.
\end{enumerate}
\end{lemma}

\begin{lemma}[Entries of the matrix exponential]\label{lem:modes}
Let $A$ be an $m \times m$ matrix with distinct eigenvalues $\lambda_1, \dots, \lambda_r \in \C$,
where $\lambda_\ell$ has algebraic multiplicity $m_\ell$. Then there are constant matrices
$C_{\ell,j} \in \C^{m \times m}$ such that, for all real $t$,
\begin{equation}\label{eq:modes}
  e^{At} = \sum_{\ell=1}^{r} \sum_{j=0}^{m_\ell - 1} t^{j} e^{\lambda_\ell t}\, C_{\ell,j} .
\end{equation}
In words: every entry of $e^{At}$ is a finite sum of terms $c\, t^{j} e^{\lambda t}$, with
$\lambda$ an eigenvalue of $A$ and $j$ smaller than the multiplicity of $\lambda$.
\end{lemma}

\begin{proof}
For $\lambda \in \C$ let
\[
  G_\lambda = \{\, v \in \C^{m} : (A - \lambda I)^{k} v = 0 \text{ for some } k \ge 0 \,\}
\]
be the \emph{generalized eigenspace} of $\lambda$. It is a subspace of $\C^{m}$ that $A$ maps into
itself, and it is not $\{0\}$ exactly when $\lambda$ is an eigenvalue of $A$. We use two standard
facts about generalized eigenspaces
(see, e.g., \cite[Chapter~8]{Axler}):
\begin{enumerate}[label=\textup{(\roman*)}]
  \item every vector of $\C^{m}$ is a sum $v_1 + \dots + v_r$ with $v_\ell \in G_{\lambda_\ell}$
        for each $\ell$;
  \item the dimension of $G_{\lambda_\ell}$ is the algebraic multiplicity $m_\ell$, that is, the
        multiplicity of $\lambda_\ell$ as a root of the characteristic polynomial of $A$.
\end{enumerate}

\emph{The power $m_\ell$ of $A - \lambda_\ell I$ vanishes on $G_{\lambda_\ell}$.} The matrix
$N = A - \lambda_\ell I$ maps $G_{\lambda_\ell}$ into itself, and some power of $N$ vanishes on
$G_{\lambda_\ell}$: take a power that kills every vector of a basis. So the restriction of $N$ to
$G_{\lambda_\ell}$ is nilpotent. Its characteristic polynomial is therefore $z^{d}$, where $d$ is
the dimension of $G_{\lambda_\ell}$, and by the Cayley--Hamilton theorem $N^{d}$ vanishes on
$G_{\lambda_\ell}$. By (ii), $d = m_\ell$. Hence
\begin{equation}\tag{$\ast$}\label{eq:nilpotent}
  (A - \lambda_\ell I)^{m_\ell}\, v = 0 \qquad \text{for every } v \in G_{\lambda_\ell} .
\end{equation}

\emph{The exponential on $G_{\lambda_\ell}$.} Let $v \in G_{\lambda_\ell}$. By
Lemma~\ref{lem:exp}\ref{lem:exp:shift}, $e^{At} = e^{\lambda_\ell t}\, e^{(A - \lambda_\ell I)t}$.
Apply the series for $e^{(A - \lambda_\ell I)t}$ to $v$; by \eqref{eq:nilpotent} the terms with
$j \ge m_\ell$ vanish, so
\[
  e^{At} v = e^{\lambda_\ell t} \sum_{j=0}^{m_\ell - 1} \frac{t^{j}}{j!}\, (A - \lambda_\ell I)^{j} v .
\]

\emph{The columns of $e^{At}$.} By (i), write the $b$th standard basis vector as
$e_b = v_{b,1} + \dots + v_{b,r}$ with $v_{b,\ell} \in G_{\lambda_\ell}$. The $b$th column of
$e^{At}$ is
\[
  e^{At} e_b = \sum_{\ell=1}^{r} \sum_{j=0}^{m_\ell - 1} t^{j} e^{\lambda_\ell t}\,
  \frac{1}{j!}\, (A - \lambda_\ell I)^{j} v_{b,\ell} .
\]
This is \eqref{eq:modes}, with $C_{\ell,j}$ the matrix whose $b$th column is
$\frac{1}{j!} (A - \lambda_\ell I)^{j} v_{b,\ell}$, for $b = 1, \dots, m$.
\end{proof}

For example, $A = \begin{bmatrix} 1 & 1 \\ 0 & 1 \end{bmatrix}$ gives
$e^{At} = \begin{bmatrix} e^{t} & t e^{t} \\ 0 & e^{t} \end{bmatrix}$. Powers of $t$ occur only
when a repeated eigenvalue has too few independent eigenvectors; if $A$ is diagonalizable, only
$j = 0$ occurs.

\begin{lemma}[Perron--Frobenius for Metzler matrices]\label{lem:PF}
Let $B$ be a Metzler matrix and let $\mu$ be the largest real part of an eigenvalue of $B$. Then
$\mu$ is an eigenvalue of $B$, and every eigenvalue $\lambda \ne \mu$ of $B$ satisfies
$\Re\lambda < \mu$. The eigenvalue $\mu$ itself may be repeated, and no irreducibility is assumed.
\end{lemma}

\begin{proof}
For real vectors $u$ and $v$ we write $u \ge v$ if $u - v \ge 0$, and $u \le 0$ if $-u \ge 0$. For
a complex vector $z$, $|z|$ is the real vector with entries $|z_i|$.

\emph{\textup{(i)} If $A$ is a Metzler matrix, then $e^{At} \ge 0$ for $t \ge 0$.} Choose
$c \ge 0$ with $P = A + cI \ge 0$. By Lemma~\ref{lem:exp}\ref{lem:exp:shift},
$e^{At} = e^{-ct} e^{Pt} = e^{-ct} \sum_{k=0}^{\infty} \frac{t^{k}}{k!} P^{k}$, and every power
$P^{k}$ is a nonnegative matrix. (Step~1 of Section~\ref{sec:proof} repeats this argument and
refines it.)

\emph{\textup{(ii)} Inverse positivity.} Let $A$ be a Metzler matrix all of whose eigenvalues have
negative real part, and let $d$ be a real vector with $Ad \le 0$. Then $d \ge 0$.

To see this, put $\psi(t) = e^{At} d$. Since $A$ commutes with $e^{At}$,
Lemma~\ref{lem:exp}\ref{lem:exp:ode} gives $\dot{\psi}(t) = e^{At} A d$, and by (i) this vector is
$\le 0$ for $t \ge 0$. So every entry of $\psi$ is nonincreasing on $[0, \infty)$. By
Lemma~\ref{lem:modes}, every entry of $e^{At}$ is a sum of terms $a\, t^{j} e^{\lambda t}$ with
$\Re\lambda < 0$. Each such term tends to $0$ as $t \to \infty$, because its modulus is
$|a|\, t^{j} e^{(\Re\lambda) t}$ and exponential decay beats any power of $t$. Hence
$\psi(t) \to 0$. A nonincreasing function on $[0, \infty)$ that tends to $0$ is nonnegative, so
$d = \psi(0) \ge 0$.

\emph{\textup{(iii)} A comparison.} Let $\sigma$ be a real number with $\Re\lambda \le \sigma$ for
every eigenvalue $\lambda$ of $B$. Suppose that $Bw \ge r w$ for a real number $r$ and a real
vector $w \ge 0$ with $w \ne 0$. Then $r \le \sigma$.

Suppose instead that $r > \sigma$. The matrix $A = B - rI$ is Metzler, and its eigenvalues are the
numbers $\lambda - r$ with $\lambda$ an eigenvalue of $B$, so they have negative real part. Also
$A(-w) = -(Bw - rw) \le 0$. By (ii), $-w \ge 0$. With $w \ge 0$ this gives $w = 0$, a
contradiction.

\emph{\textup{(iv)} Moduli of an eigenvector.} Let $Bz = \lambda z$ with $\lambda \in \C$ and $z$
a complex vector, and choose $c \ge 0$ with $P = B + cI \ge 0$. Then $Pz = (\lambda + c) z$, so for
every index $i$
\[
  |\lambda + c|\, |z_i| = \Bigl| \sum_{j} P_{ij} z_j \Bigr| \le \sum_{j} P_{ij}\, |z_j|
  = (B|z|)_i + c\, |z_i| .
\]
That is, $B|z| \ge (|\lambda + c| - c)\, |z|$.

\emph{\textup{(v)} Conclusion.} Let $\lambda$ be an eigenvalue of $B$ with $\Re\lambda = \mu$, and
let $z \ne 0$ be an eigenvector for it. The vector $|z| \ge 0$ is not zero, so (iv) and (iii), with
$\sigma = \mu$, $w = |z|$ and $r = |\lambda + c| - c$, give $|\lambda + c| - c \le \mu$. Hence
\[
  |\lambda + c| \le \mu + c = \Re(\lambda + c) .
\]
A complex number whose modulus is at most its real part is real. So $\lambda$ is real, that is,
$\lambda = \mu$. By the definition of $\mu$, some eigenvalue of $B$ has real part $\mu$; we have
just seen that it equals $\mu$. So $\mu$ is an eigenvalue of $B$, and it is the only eigenvalue
with real part $\mu$: every other eigenvalue has real part less than $\mu$.
\end{proof}

Lemma~\ref{lem:PF} is the form, for Metzler matrices, of the classical theorem that the spectral
radius of a nonnegative matrix is an eigenvalue of that matrix
\cite[Theorem~8.3.1]{HJ}. It also follows from that theorem, applied to $P = B + cI$.

% =====================================================================
\section{Proof of Theorem~\ref{thm:main}}\label{sec:proof}
% =====================================================================

\subsection*{Step 1. $e^{Mt}$ is nonnegative, and reachability decides which entries are positive}

Choose $c \ge 0$ with $P = M + cI \ge 0$; for instance $c = \max_i |M_{ii}|$. By
Lemma~\ref{lem:exp}\ref{lem:exp:shift},
\[
  e^{Mt} = e^{-ct} e^{Pt} = e^{-ct} \sum_{k=0}^{\infty} \frac{t^{k}}{k!} P^{k} .
\]
Every power $P^k$ is a nonnegative matrix, so $e^{Mt} \ge 0$ for $t \ge 0$. Moreover, for $t > 0$
the $(i,j)$ entry of $e^{Mt}$ is positive if $[P^k]_{ij} > 0$ for at least one $k \ge 0$, and it is
zero for all $t \ge 0$ if $[P^k]_{ij} = 0$ for every $k \ge 0$.

For $k \ge 1$, the entry $[P^k]_{ij}$ is the sum, over all sequences of indices $j = i_0, i_1, \dots, i_k = i$,
of the products
\[
  P_{i_k i_{k-1}} \cdots P_{i_2 i_1} P_{i_1 i_0} .
\]
These products are nonnegative, so the sum is positive exactly when some sequence has every factor
positive. In such a sequence each step
either stays put ($i_s = i_{s-1}$, a diagonal entry of $P$) or follows an arc
($i_s \ne i_{s-1}$ and $M_{i_s i_{s-1}} > 0$, since $P$ and $M$ have the same off-diagonal
entries). Deleting the steps that stay put shows that $i$ is reachable from $j$. Conversely, if $i$
is reachable from $j$ along $k \ge 1$ arcs then $[P^k]_{ij} > 0$, and the case $i = j$ is covered by
$P^0 = I$. Therefore, for every $t > 0$,
\begin{equation}\label{eq:pattern}
  [e^{Mt}]_{ij} > 0 \ \text{ if $i$ is reachable from $j$,}
  \qquad
  [e^{Mt}]_{ij} = 0 \ \text{ otherwise.}
\end{equation}

\subsection*{Step 2. Only the reachable indices matter}

Since $x(t) = e^{Mt} x(0)$ and $x_j(0) = 0$ for $j \notin S$,
\[
  x_i(t) = \sum_{j \in S} [e^{Mt}]_{ij}\, x_j(0) ,
\]
a sum of nonnegative terms. If $i \in R$, some $j \in S$ has $i$ reachable from $j$, and by
\eqref{eq:pattern} the corresponding term is positive for $t > 0$. If $i \notin R$, every term is
zero. Hence, for every $t > 0$,
\begin{equation}\label{eq:support}
  x_i(t) > 0 \ \text{ if } i \in R,
  \qquad
  x_i(t) = 0 \ \text{ if } i \notin R,
\end{equation}
and $x_i(t) = 0$ for all $t \ge 0$ when $i \notin R$. Three consequences follow.

\begin{itemize}[leftmargin=2em]
  \item \emph{The total is positive.} $W(0) > 0$ because $x(0) \ge 0$ is nonzero, and $W(t) > 0$
        for $t > 0$ by \eqref{eq:support}. This proves
        part~\ref{thm:main:support} of the theorem.
  \item \emph{No arc leaves $R$.} If $j \in R$ and $i \notin R$, then
        \begin{equation}\label{eq:noarc}
          M_{ij} = 0 ,
        \end{equation}
        because otherwise $i$ would be reachable from $j$, and hence from $S$.
  \item \emph{The reduced system.} Put $y(t) = x_R(t)$. For $i \in R$,
        $\dot{x}_i = \sum_{j} M_{ij} x_j = \sum_{j \in R} M_{ij} x_j$, because $x_j \equiv 0$ for
        $j \notin R$. Hence
        \[
          \dot{y} = By, \qquad y(t) = e^{Bt} y(0), \qquad y(t) \gg 0 \ \text{ for every } t > 0 .
        \]
        The matrix $B$ is Metzler, because its off-diagonal entries are off-diagonal entries of $M$.
\end{itemize}

\subsection*{Step 3. The leading term of $e^{Bt}$}

Apply Lemma~\ref{lem:PF} to $B$. This proves part~\ref{thm:main:mu} of the theorem, and shows that
every eigenvalue $\lambda \ne \mu$ of $B$ has $\Re\lambda < \mu$. Apply Lemma~\ref{lem:modes} to
$B$ and write $C_j$ for the coefficient matrix of $t^{j} e^{\mu t}$:
\begin{equation}\label{eq:expansion}
  e^{Bt} = \sum_{j} t^{j} e^{\mu t}\, C_j
  \;+\; \sum_{\lambda \ne \mu} \sum_{j} t^{j} e^{\lambda t}\, C_{\lambda,j} .
\end{equation}

\emph{Some $C_j$ is nonzero.} Take $u \ne 0$ with $Bu = \mu u$. By
Lemma~\ref{lem:exp}\ref{lem:exp:eig}, $e^{-\mu t} e^{Bt} u = u$ for all $t$. If every $C_j$ were
zero, then by \eqref{eq:expansion} each entry of $e^{-\mu t} e^{Bt}$ would be a sum of terms
$c\, t^{j} e^{(\lambda - \mu)t}$ with $\Re\lambda < \mu$. Each such term tends to $0$ as
$t \to \infty$, because its modulus is $|c|\, t^{j} e^{-(\mu - \Re\lambda)t}$ and exponential decay
beats any power of $t$. This would give $u = 0$, a contradiction.

\emph{Definition of $k$ and $L$.} Let $k$ be the largest $j$ with $C_j \ne 0$, and put $L = C_k$.
Divide \eqref{eq:expansion} by $t^{k} e^{\mu t}$. Every term other than $L$ carries a factor
$t^{j-k}$ with $j < k$, or a factor $t^{j-k} e^{(\lambda - \mu)t}$ with $\Re\lambda < \mu$; both
tend to $0$. Therefore
\begin{equation}\label{eq:L}
  \lim_{t \to \infty} \frac{e^{Bt}}{t^{k} e^{\mu t}} = L \ne 0 .
\end{equation}
The fastest-growing terms of $e^{Bt}$ grow like $t^{k} e^{\mu t}$, and $L$ holds their
coefficients. Both $k$ and $L$ depend only on $B$.

\subsection*{Step 4. $L$ is nonnegative, and its nonzero columns are eigenvectors}

\emph{\textup{(i)} $L \ge 0$.} By Step~1 applied to the Metzler matrix $B$, $e^{Bt} \ge 0$; and
$t^{k} e^{\mu t} > 0$ for $t > 0$. A limit of nonnegative matrices is nonnegative; in particular
$L$ is real.

\emph{\textup{(ii)} $BL = \mu L$.} Fix $s > 0$. Since $e^{Bs} e^{Bt} = e^{B(s+t)}$ by
Lemma~\ref{lem:exp}\ref{lem:exp:semigroup},
\[
  e^{Bs} \cdot \frac{e^{Bt}}{t^{k} e^{\mu t}}
  \;=\;
  \frac{(s+t)^{k} e^{\mu(s+t)}}{t^{k} e^{\mu t}}
  \cdot
  \frac{e^{B(s+t)}}{(s+t)^{k} e^{\mu(s+t)}} .
\]
Let $t \to \infty$. The left side tends to $e^{Bs} L$ by \eqref{eq:L}. On the right, the first
factor equals $(1 + s/t)^{k} e^{\mu s}$ and tends to $e^{\mu s}$, and the second factor tends to
$L$ by \eqref{eq:L}. Hence
\begin{equation}\label{eq:semigroup}
  e^{Bs} L = e^{\mu s} L \qquad \text{for every } s > 0 .
\end{equation}
Subtract $L$ from both sides and divide by $s$:
\[
  \frac{e^{Bs} - I}{s}\, L = \frac{e^{\mu s} - 1}{s}\, L .
\]
As $s \to 0^{+}$ the left side tends to $BL$, because
$\frac{d}{ds} e^{Bs} \big|_{s=0} = B$ by Lemma~\ref{lem:exp}\ref{lem:exp:ode}, and the right side
tends to $\mu L$. Therefore
\begin{equation}\label{eq:eig}
  BL = \mu L .
\end{equation}
Consequently, for every vector $z$ we have $B(Lz) = \mu\,(Lz)$: the vector $Lz$ is either zero or
an eigenvector of $B$ for the eigenvalue $\mu$.

\emph{\textup{(iii)} If $z \gg 0$, then $Lz \ne 0$.} By \eqref{eq:L} and (i), some entry $L_{ij}$
is positive. Then $(Lz)_i = \sum_{l} L_{il} z_l \ge L_{ij} z_j > 0$.

\subsection*{Step 5. The limit for the reduced system}

Fix any $\tau > 0$ and put $z = y(\tau)$; by Step~2, $z \gg 0$. For $t > 0$,
Lemma~\ref{lem:exp}\ref{lem:exp:semigroup} gives $y(\tau + t) = e^{Bt} z$, so by \eqref{eq:L}
\[
  \frac{y(\tau + t)}{t^{k} e^{\mu t}} \;\longrightarrow\; v := Lz
  \qquad (t \to \infty).
\]
By Step~4, $v \ge 0$, $v \ne 0$ and $Bv = \mu v$. Summing the entries,
$\one^\T y(\tau + t) / (t^{k} e^{\mu t}) \to \one^\T v > 0$. Dividing the two limits, which is
legitimate because the limit of the denominator is not zero,
\begin{equation}\label{eq:q}
  \frac{y(t)}{\one^\T y(t)} \;\longrightarrow\; q := \frac{v}{\one^\T v}
  \qquad (t \to \infty),
  \qquad\text{with}\quad q \ge 0, \quad \one^\T q = 1, \quad Bq = \mu q .
\end{equation}

\subsection*{Step 6. Back to the full system}

By \eqref{eq:support}, $W(t) = \one^\T y(t)$. Hence $p_R(t) = y(t)/\one^\T y(t)$, and $p_i(t) = 0$
for $i \notin R$. By \eqref{eq:q}, $p(t) \to p^*$, where $p^*_R = q$ and $p^*_i = 0$ for
$i \notin R$. So $p^* \ge 0$ and $\one^\T p^* = 1$.

It remains to check $M p^* = \mu p^*$. For any $i$, $(M p^*)_i = \sum_{j \in R} M_{ij} q_j$.
If $i \in R$, this is $(Bq)_i = \mu q_i = \mu p^*_i$. If $i \notin R$, every $M_{ij}$ in the sum is
zero by \eqref{eq:noarc}, so $(M p^*)_i = 0 = \mu p^*_i$. This proves
part~\ref{thm:main:limit}. \qed

\begin{remark}[What each hypothesis does]\label{rem:roles}
The Metzler property is used in four places. It gives $e^{Mt} \ge 0$ and, because sums of
nonnegative products cannot cancel, the exact positivity pattern \eqref{eq:pattern} (Step~1). It
makes the absence of an arc from $j$ to $i$ mean $M_{ij} = 0$, which is \eqref{eq:noarc} (Step~2).
It makes the eigenvalue of largest real part real and strictly dominant, so that the leading term
of $e^{Bt}$ does not oscillate (Lemma~\ref{lem:PF}, Step~3). And it gives $L \ge 0$ (Step~4). Strict positivity of
$y(\tau)$ is used once, in Step~5, to guarantee $L y(\tau) \ne 0$. In a general linear system both
can fail: nonreal eigenvalues cause persistent oscillation in the direction of the solution when
they contribute to its leading term, and an initial state can have no component along the leading
term of the matrix exponential. The reduction to $R$ is what supplies strict positivity. Applied
to the full matrix $M$, the argument would fail: in Example~\ref{ex:reduction} the leading
coefficient matrix of $e^{Mt}$ annihilates $x(0)$, and the solution grows at a rate smaller than
the largest eigenvalue of $M$.
\end{remark}

\begin{example}[The reduction to $R$ is needed]\label{ex:reduction}
Let $M = \begin{bmatrix} 1 & 0 \\ 0 & 2 \end{bmatrix}$. For $x(0) = (1, 0)^\T$ we get $R = \{1\}$,
$B = [\,1\,]$, $\mu = 1$ and $p(t) = e_1$ for all $t$, although the dominant eigenvalue of $M$ is
$2$. For $x(0) \gg 0$ we get $R = \{1,2\}$, $\mu = 2$ and $p^* = e_2$.
\end{example}


% =====================================================================
\begin{thebibliography}{9}

\bibitem{Axler}
S.~Axler,
\emph{Linear Algebra Done Right}, 4th ed.,
Undergraduate Texts in Mathematics, Springer, 2024.

\bibitem{HJ}
R.~A. Horn and C.~R. Johnson,
\emph{Matrix Analysis}, 2nd ed.,
Cambridge University Press, Cambridge, 2013.

\end{thebibliography}

\end{document}
